\section{From the integral lattice to the residual local model}

% BEGIN M80R1
\paragraph{M80R1. Full residual lattices and their Frobenius bounds.}
Use the level-one fields of W80N1 and a residual presentation of
$S.\mathrm{RealD}$ with $v$ as invariant parameter. Write $v$ also
for its image in $k[[v]]$. Let $T$ be crystalline and let $q$ identify
$T/\lambda T$ equivariantly with a residual representation $A_0$.
For each letter take $M_i$ to be the image of the invariant Wach
lattice by the specified presentation map $j_i$. By W80N2 this
image is spanned over $k[[v]]$ by the images of one invariant basis
$b_j$. Indeed all coefficients in $O[[v]]$ reduce onto $k[[v]]$.
These vectors are a basis over $k((v))$: in the full module the $b_j$
extend to an $\mathcal E_O$-basis. The kernel rule of the presentation,
applied to multiplication by $\lambda$ on $T$, says exactly that
$j_i$ kills an invariant vector if and only if its letter component
lies in $\lambda D_i(T)$. Thus any integral relation among the reduced
columns, lifted coefficientwise to $O[[v]]$, has all its coefficients
zero modulo $\lambda$ in $k((X))$. The scalar map $k[[v]]\to k[[X]]$
is injective by W80N1. Clearing a power of $v$ handles any Laurent
relation. The number of columns is $\dim_k A_0$, so the module is
full. This verifies the exact image prescription in $\mathrm{RedN}$.

The columns $\varphi b_j$ likewise form a basis of $P_i$ and are
invariant. For the inclusions $Q^{M}P_i\subset N_i$ and
$Q^hN_i\subset P_i$ from W80N2 replace $Q$ by its associate $v+p$.
The new change-of-basis coefficients lie in $O[[v]]$, since both
bases and $v+p$ are invariant and the bases are independent.
Reducing the inclusions via $j_i$, which commutes with Frobenius,
gives
\[
 v^h M_i\subseteq k[[v]]\varphi M_{i-1}\subseteq v^{-M}M_i.
\]
The same argument in rank one reduces the equality of W80N2 to
$k[[v]]\varphi M_{i-1}=v^{a_i}M_i$, for either sign of $a_i$.
The divided action $\nu$ reduces to $\nabla=(\gamma-1)/v$ by the
scalar and action rules for $j_i$. Thus $(\gamma-1)M_i\subset vM_i$.
These are the three native lattice inequalities.
% END M80R1

% BEGIN M80R2
\paragraph{M80R2. Residual Sen multiplicity and the boundary polynomial.}
In $k[[X]]$ one has $\gamma X=X+X^p+X^{p+1}$. Consequently
$\gamma v-v\in (v^2)$: substitution raises the order of the
difference on $X^{p-1}k[[X]]$ to at least $2(p-1)$, and both series
are invariant. It follows that $\nabla$ preserves $vM_i$ and induces
a $k$-linear operator on $M_i/vM_i$. Its basis matrix is the reduction
of $B(-p)$ in W80N5, since $v+p$ reduces to $v$. For every integer
$j$, $(1+p)^j\equiv1+jp\bmod p^2$, also for $j<0$ by inversion.
Hence $c(-u)\bmod\lambda=u\bmod p$. The dimension of the generalized
$x$-eigenspace on $M_i/vM_i$ equals the number of jumps congruent
to $x$, with multiplicity, by the split characteristic polynomial.
The source's quotient defining $\mathrm{sen}$ is just that generalized
eigenspace: a class represented by $m$ belongs precisely when a power
of $\nabla-x$ carries $m$ into $vM_i$.

For $s\ge0$ form $I=\{a-u:0\le a<s,\ u\in U_i\}$ as a set.
By W80N4 the operator $\prod_{m\in I}(\nabla+\bar m)$ carries
$M_i$ into $v^sM_i$. In reducing the integral output, a quotient by
$(v+p)^s$ on the letter is invariant: its product is invariant and
the scalar is regular in the free Wach module. Suppose the only
multiples of $p$ in $I$ are zero. Remove $m=0$ from the product to
form $f(Z)$, so $f(0)\ne0$. If $0\in I$, $\nabla f(\nabla)$ is
the original operator. Otherwise apply the original bound to
$\nabla y\in M_i$. In either case $\nabla f(\nabla)M_i\subseteq
v^s M_i$, exactly $\mathrm{EPoly}$, including the empty product.
% END M80R2

% BEGIN M80R3
\paragraph{M80R3. Crystalline calculus and reductions of full summands.}
On a split unramified good layer the four $S.\mathrm{CrysCalc}$
rules follow from period bases. A product of two crystalline
representations has the union of their invariant bases and the
direct sum filtration. Tensoring by a crystalline character with
labelled jump $a_i$ multiplies the periods on each component by its
invariant line generator. The de Rham order of this generator is
$a_i$, so graded dimensions shift by $a_i$. Its inverse gives the
opposite shift. The cyclotomic line has invariant generator $t^{-1}$
times its representation generator, hence jump $-1$, with valuation
one for $t$. The unramified line has jump zero by the integral
unramified period basis in S80J4. These calculations use the tied
period inclusions and are preserved by rational equivariant isomorphism.

Let $T$ be a sum of $\mu$ full objects, in the sense of the native
$\mathrm{SumFull}$, with an arbitrary choice of stable lattice in
the resulting rational product. Its reduction has the same
composition factors as the reduction of the direct-sum lattice:
scale one lattice into the other, with finite cokernel $C$. The
kernel and cokernel of multiplication by $\lambda$ on $C$ have
the same composition factors, by the two exact sequences through
$\lambda C$. Reducing the lattice inclusion and using its exact
kernel $C[\lambda]$ and cokernel $C/\lambda C$ gives the claim.
The simple factors of these finite modules are killed by $\lambda$,
as follows from nilpotence of $\lambda$ and simplicity. The full
summands have the prescribed reference factors by definition.
Twisting the entire lattice by $\chi_1^{-1}$ shifts its jumps to
$\mu(U_{t,i}-w_i)$ and its reduction to the factors of $\mu$ copies
of the reference representation with rank block $\overline{T_0/\chi_1}$
and lines $\omega^{-e}$, $0\le e<t$, where $\omega=\bar\varepsilon$.
Here the multiset notation means repeating each shifted jump $\mu$
times. W80N3 and naturality of $j_i$ give a direct sum of the
residual lattices for an integral direct sum: the inclusions and
the coordinate projections commute with both actions and reduce
to the corresponding residual maps. Thus these base pairs are
closed under sums, with multiplicities adding.
% END M80R3

% BEGIN M80R4
\paragraph{M80R4. The standard pair, all character pairs and the realization of admitted quotients.}
Use as base pairs precisely M80R3's twisted reductions with their
M80R1 lattices. For a jump $u\in U_{t,i}$ one has
$1\le u\le h_i-1$ and $1\le w_i\le h_i-1$. Thus the shifted
jumps lie in $[-w_i,h_i-w_i]$. M80R1 proves the required base
bounds. M80R2 and M80R3 give the stipulated Sen counts and
semisimplifications. At a boundary letter put $s=h_i-w_i+1$.
For $0\le a<s$ and $u$ in the low list or the character string,
$u,w_i$ lie in $[h_i-\eta_i,h_i-1]$ and $w_i+a\in[w_i,h_i]$.
Thus $|w_i+a-u|\le\eta_i<p$ by T1, and divisibility by $p$ implies
zero. For a cloud entry the same divisibility would give
$h_i-u\equiv h_i-w_i-a\bmod p$, with the right side in
$[0,\eta_i]$, excluded by T3. M80R2 now supplies the boundary
polynomial with the prescribed $s$.

Reduce the twisted standard rank-$r$ lattice and, for each $0\le e<t$,
the line $\chi_{e+1}/\chi_1=\varepsilon^{-e}$. Their integral direct
sum before twisting is full: its crystalline jumps are $U_{t,i}$,
and its reduction on the standard lattice is the reference
representation by definition. For any other stable lattice use
M80R3's composition-factor argument. Each indicated direct summand
therefore has a model by inclusion and projection on this base pair.
The standard lattice has precisely the $U_i-w_i$ Sen counts, and
the character lattices have jump $e$, Sen count at $\bar e$, and
Frobenius shift $e$ by M80R1--M80R2. Use the cyclotomic line itself
for $e=-1$ (the model requirement begins at zero). For $e=0$ write
the rank-one lattice as $v^{d_i}k[[v]]$ times the distinguished
fixed unit. Shift zero forces $d_i=p d_{i-1}$, since Frobenius
raises the order by $p$. Going round the finite cycle yields
$d_i=0$. Outside the specified range the unused lattice entries
may be chosen arbitrarily. The inertia surjectivity argument of
S80J5 applies to any unramified $K$, so $\omega$ has order $p-1$,
as required in the stated range of powers.

Finally an admitted residual module, after its specified twist,
comes with an equivariant quotient from the integral lattice of a
sum of full objects. Its target is killed by $\lambda$, so the map
factors onto it through that very twisted full reduction. The
factor is $k$-linear since $O\to k$ is onto. Push the base lattice
forward by this surjection and take the identity for the inclusion
in the subquotient definition of model. This also retains the
standard representation's required reduction with kernel $\lambda T_0$.
All these choices use a single residual functor and the same $j_i$
for the given $k$. Hence for every native residual presentation the
constructed record is $\mathrm{ModTh}$ with the additional clauses of
$S.\mathrm{Real}$, and $S.\mathrm{RealD}$ implies $S.\mathrm{Real}$
on the compatible level-one system.
% END M80R4
